\CharacterSheet{ghoul}{Carniçal}{

\shortdescription{O bobo e animal de estimação da família, o melhor artista.}

\paragraph{Objetivos}
\begin{itemize}
  \item Descobrir quem matou o recém-chegado.
\end{itemize}

\paragraph{Condições de derrota}
Não conseguir expulsar por votação quem matou \Character{newcomer}, caso exista uma assassina. Não perdes se \Character{enforcer} for expulso por votação, mesmo que não seja quem matou \Character{newcomer}.

\paragraph{Segredos}
\begin{itemize}
  \item ?.
\end{itemize}

\paragraph{Ações}
\begin{itemize}
  \item Sê estranho e engraçado.
\end{itemize}

\paragraph{Contexto}

Fui em tempos um encantador bobo ambulante. \timestamp{No ano 100 d.C.}, \Character{enforcer} apaixonou-se por mim e mandou-me trazer para esta família. Pensei que tornar-me vampiro me faria livre. Em vez disso, descobri que tenho de obedecer às ordens faladas de \Character{enforcer}. A família trata-me como um animal de estimação e uma piada, não como um deles, e tive séculos para odiar isso.

\timestamp{Ontem, à 01:00}, quando \Character{usurper} deixou o seu gabinete aberto, fui procurar uma forma de quebrar o meu vínculo. Roubei um ritual chamado ``Mudar o futuro''. Vi \Character{prevlover} roubar outro ritual, para mudar e conhecer uma pessoa. \timestamp{Ontem, às 15:00}, quando realizei o meu ritual, vi que \Character{newcomer} iria morrer.

\Character{enforcer} já me tinha ordenado que revelasse toda a informação que conhecia sobre a família, por isso não tive escolha senão contar-lhe o que vi. Ainda espero que o ritual que roubei me possa levar à liberdade, mesmo que até agora só me tenha mostrado morte.

\GMNotes{\Character{ghoul} fornece o aviso que alerta primeiro \Character{enforcer} para a morte de \Character{newcomer} e é a única testemunha do roubo de \Character{prevlover}. O seu ressentimento e o ritual roubado tornam-no suspeito, mas não causou a morte. O vínculo é importante: quando \Character{enforcer} dá uma ordem clara, \Character{ghoul} tem de obedecer; o MJ pode usar essa limitação para libertar informação sem forçar o jogador a sair da personagem.}

}
